%! Solving Radical Equations \& Extraneous Solutions — Unit 6, Lesson 6.5 — Exit Ticket (Answer Key)
\documentclass[10pt]{article}
\usepackage{algebra2-article}
\usepackage{algebra2-key}   % pulls in -boxes; do NOT also load -boxes
\newcommand{\TallMath}[1]{$\displaystyle #1\rule[-1.4em]{0pt}{3.2em}$}
% In-formula answer slot (\ans is text-mode and must never appear inside $...$).
\newcommand{\mans}[1]{{\color{keyred}\mathbf{#1}}}

\begin{document}

\pageheader{Unit 6, Lesson 6.5}{Exit Ticket --- Answer Key}
\namedateperiod

\vspace{0.08in}

No notes. Show the check for every candidate --- an answer with no check is not finished.

\begin{enumerate}[leftmargin=1.8em, itemsep=10pt]
  \item \textbf{Solve and check.} \; $\sqrt{3x+1}-2=3$

        {\small \emph{Work:} isolate: $\sqrt{3x+1}=5$. \; Square: $3x+1=25$. \; Solve: $3x=24$.}

        $x=\mans{8}$ \hspace{0.3in} Check: \ans{$\sqrt{3(8)+1}-2=\sqrt{25}-2=5-2=3$ \checkmark}

  \item \textbf{Solve, check, and sort.} \; $\sqrt{x+5}=x-1$

        {\small \emph{Work:} square: $x+5=x^{2}-2x+1$, so $0=x^{2}-3x-4=(x-4)(x+1)$. \\
        Check $x=4$: $\sqrt{9}=3$ and $4-1=3$. \checkmark \quad
        Check $x=-1$: $\sqrt{4}=2$ but $-1-1=-2$. \; $2\ne-2$.}

        (a) Candidates: \ans{$4$ and $-1$} \hspace{0.2in} (b) Solution set: \ans{$\{4\}$}

        \vspace{5pt}
        (c) One candidate was extraneous. In one sentence, say what went wrong with it --- your
        sentence should mention a \emph{sign}. \ansline{At $x=-1$ the right side is $-2$ while the
        left side is the principal root $\sqrt4=2$; squaring turned $2$ and $-2$ into the same
        number, so the squared equation accepted a value the original never did.}

  \item \textbf{SOL-style multiple choice.} Which equation has \textbf{no real solution}?
        \begin{enumerate}[label=\Alph*., leftmargin=1.9em, itemsep=1pt]
          \item $\sqrt{x-1}=5$
          \item $\sqrt[3]{x-1}=-5$
          \item $\sqrt{x-1}=-5$ \quad \textcolor{keyred}{\textbf{$\leftarrow$ correct}}
          \item $\sqrt{x-1}=0$
        \end{enumerate}
        \vspace{2pt}
        \begin{itemize}[leftmargin=1.4em, nosep]
          \item \textbf{C} is correct: a principal \emph{even} root is never negative, so nothing
                can make $\sqrt{x-1}$ equal $-5$. (Squaring anyway gives the extraneous $x=26$.)
          \item \textbf{A} solves normally: $x-1=25$, $x=26$, and it checks.
          \item \textbf{B} is the discriminating distractor --- the index is \emph{odd}, so a
                negative value is perfectly legal: $x-1=-125$, $x=-124$. A student who has learned
                ``roots can't be negative'' instead of ``\emph{even} roots can't'' picks this.
          \item \textbf{D} gives $x=1$, the endpoint of the graph.
        \end{itemize}
\end{enumerate}

\begin{teachernote}
\small Graded for completion. Sort into four piles; they decide how 6.6 opens and what needs a
re-teach first.
\begin{itemize}[leftmargin=1.3em, nosep]
  \item \textbf{Got it} --- item 1 $x=8$, item 2 $\{4\}$ with $-1$ named extraneous, item 3 $=$ C.
  \item \textbf{Skipped the check} --- item 2 answered $\{4,-1\}$. The most common miss, and the
        cheapest to fix: the check is a \emph{step}, not a courtesy. Have them substitute $-1$ out
        loud.
  \item \textbf{Rejected on reflex} --- item 2 answered $\{4\}$ but with the reason ``$-1$ is
        negative.'' Being negative is not what makes a candidate fail; making the \emph{other side}
        negative is. Homework item 2(b) is built to break this habit --- both of its candidates are
        negative and both are valid.
  \item \textbf{Index confusion} --- item 3 $=$ B. This is the Unit 6 spine failing, not a
        procedural slip. Re-run the Section 4 contrast ($\sqrt{x-4}=-2$ beside
        $\sqrt[3]{x-4}=-2$) before moving on.
  \end{itemize}
\textbf{Item 2(c) is the conceptual item.} A student who identifies the extraneous candidate but
cannot name the sign mismatch has learned to check without learning why checking is necessary --- and
will drop it again the first time the problem gets long.
\end{teachernote}

\end{document}
